Inserir o texto aqui.

\section{Exemplo de Tabela}
\label{sec:exemplo_tabela}
Apenas um exemplo de como criar uma tabela em \LaTeX{}. O exemplo pode ser
visto na tabela \ref{tab:exemplo_tabela}.

  \begin{table}[htbp]
   \centering
    \begin{tabular}{|c|c|}
      \hline
      \multicolumn{1}{|c|}{\textbf{Coluna 1}} &
      \multicolumn{1}{|c|}{\textbf{Coluna 2}}\\
      \hline
      bla bla bla bla & bla bla bla \\
      \hline
    \end{tabular}
    \caption{Exemplo de uma tabela}
    \label{tab:exemplo_tabela}
  \end{table}

\section{Exemplo de Figura}
Um simples exemplo de como inserir uma figura no documento. A figura
\ref{fig:uninorte} contempla este exemplo.

\begin{figure}
\centering
\includegraphics[width=8cm]{imagens/marca_uninorte.png}
\caption{Marca do UniNorte}
\label{fig:uninorte}
\end{figure}

\section{Exemplo de Refer�ncia}
\subsection{Refer�ncia externa}
Usar refer�ncia de livros, links, etc. Exemplo: De acordo com \cite{Knuth78}...

\subsection{Refer�ncia interna}
Podemos referenciar um cap�tulo, se��o ou qualquer parte do documento
utilizando o caption. Exemplo: citando o Cap�tulo \ref{cap:cap_1} ou a se��o
\ref{sec:exemplo_tabela}. 